\documentclass[a4paper,12pt]{article}
\usepackage[utf8]{inputenc}
\usepackage[T1]{fontenc}
\usepackage[french]{babel}
\usepackage{geometry}
\geometry{a4paper, margin=1in}
\usepackage{hyperref}

\title{Offre de Stage : Analyse Statistique et Économétrique du Rendement des Cryptomonnaies}
\author{Pierre-O Goffard}
% \date{}

\begin{document}

\maketitle

\section*{Université de Strasbourg - Institut de Recherche Mathématiques Avancées}

\subsection*{Lieu : Université de Strasbourg}

\subsection*{Durée : 2-3 mois}
\subsection*{Début : Juin 2025}

\subsection*{Description du poste :}
Stagiaire de niveau L3 pour travailler sur l'analyse statistique des log-rendements des cryptomonnaies. Ce projet est financé par l'ANR \href{https://anr.fr/Projet-ANR-24-CE38-7885}{BLOCKFI}.

\subsection*{Missions :}

\begin{enumerate}
    \item \textbf{Collecte et Préparation des Données :}
    \begin{itemize}
        \item Récupérer des données journalières de cours des cryptomonnaies à partir de différentes sources.
        \item Nettoyer et préparer les données pour l'analyse statistique.
    \end{itemize}
    \item \textbf{Formation :}
    \begin{itemize}
        \item Se former à Git/GitHub.
        \item Se former à Python pour l'analyse de données, utiliser VSCode/Copilot/Jupyter notebooks.
        \item Se former à l'étude des séries temporelles avec Python.
    \end{itemize}

    \item \textbf{Analyse Statistique :}
    \begin{itemize}
        \item Effectuer des analyses descriptives et exploratoires des données.
        \item Identifier les tendances et les anomalies dans les log-rendements.
    \end{itemize}

    \item \textbf{Modélisation et Prédictions :}
    \begin{itemize}
        \item Développer des modèles statistiques de séries temporelles pour prédire les rendements futurs des cryptomonnaies.
        \item Évaluer la performance des modèles et ajuster les paramètres si nécessaire.
    \end{itemize}

    \item \textbf{Visualisation des Données :}
    \begin{itemize}
        \item Créer des visualisations claires et informatives pour présenter les résultats de l'analyse.
        \item Utiliser des outils comme Python (Matplotlib, Seaborn) ou Tableau pour la visualisation.
    \end{itemize}

    \item \textbf{Rapport et Présentation :}
    \begin{itemize}
        \item Rédiger un rapport détaillé des analyses effectuées et des conclusions obtenues.
        \item Préparer un support de présentation pour restituer les résultats de l'étude.
        \item Les documents devront être rédigés sur LaTeX.
    \end{itemize}
\end{enumerate}

\subsection*{Profil recherché :}

\begin{itemize}
    \item Étudiant(e) en L3 en Mathématiques, Statistiques, Informatique.
    \item Connaissances de base en statistiques et en analyse de données.
    % \item Maîtrise des outils d'analyse de données (Python, R, Excel).
    \item Capacité à travailler de manière autonome et à prendre des initiatives.
    \item Bonnes compétences en communication écrite et orale.
\end{itemize}

\subsection*{Comment postuler :}

Envoyez votre CV et une lettre de motivation à \href{mailto:goffard@unistra.fr}{goffard@unistra.fr} en précisant la référence "Stage Analyse Cryptomonnaies - ANR BLOCKFI" dans l'objet du mail.

\end{document}
